\section{Introduction}\label{sec:introduction}

Attention in real classrooms is fugitive. It rises and falls within a single lesson, drifts away from the planned task without overt behavioural signature, and recovers in ways that no end-of-unit test can resolve. Teachers integrate these fluctuations into their everyday judgement, but the instruments routinely available to characterise cognitive performance in ordinary teaching---summative tests, retrospective self-report, lesson-end ratings---operate at a temporal grain orders of magnitude coarser than the underlying construct, and they are known to carry systematic biases linked to social desirability, recall and observer expectation \parencite{Fredricks2004_school,DMello2012_dynamics,Hutt2019_mind_wandering}. Passive multimodal sensing is appealing because it is continuous, ecologically valid and non-disruptive: ambient, physiological, neurophysiological and behavioural signals can in principle be recorded during ordinary teaching with negligible interference and aligned with external attentional benchmarks for supervised modelling.

The constituent technologies of such a sensing stack are individually mature. Indoor environmental quality has been linked causally to cognitive performance in school-age populations through controlled-exposure and field-experimental designs, with carbon dioxide concentration, ventilation rate, temperature and particulate matter as the most consistently implicated variables \parencite{Wargocki2007_effects,Wargocki2020_classroom_co2,Bako2012_classroom,Satish2012_co2,Petersen2016_classroom,Allen2016_cogfx}. Dry and semi-dry electroencephalography systems now achieve signal quality comparable to conventional wet-gel setups for many spectral and event-related features, opening EEG to longitudinal field deployment outside the laboratory \parencite{LopezGordo2014_dry_eeg,DiFlumeri2019_dry_eeg_validation,Hinrichs2020_motion_artifacts}; mobile EEG has, in particular, been validated for adolescent and paediatric work \parencite{LauZhu2019_mobile_eeg_review}. Wearable physiological recorders sample heart rate, heart-rate variability and electrodermal activity at densities that previously required tethered laboratory equipment \parencite{SanchezReolid2024_eeg_fnirs}. Body-pose estimation from overhead video has demonstrated that fine-grained behavioural features---posture changes, head orientation, gross-motor activity, hand-raising---can be retrieved at scale from classroom recordings \parencite{Cao2021_openpose,Ahuja2019_edusense}. Each of these strands feeds into the broader programme of multimodal learning analytics \parencite{Blikstein2016_multimodal,Schneider2015_unraveling,Prinsloo2023_mmla_privacy,HernandezMustieles2024_wearable_education}.

\textcolor{blue}{Despite their individual maturity, these modalities have rarely been deployed in combination. No published deployment simultaneously satisfies four conjunctive properties: (a) integration of ambient classroom IoT, biometric wearables, electroencephalography and a behavioural video stream within one acquisition pipeline; (b) deployment over a longitudinal field window (months rather than session-level recording blocks) in ordinary teaching of compulsory secondary education; (c) on-site or institutional-server processing under European data-protection guarantees for minors; and (d) anchoring of the resulting signal stream to an external, independently validated attentional benchmark. The Psychomotor Vigilance Task (PVT) is well suited to serve as this external anchor, given its established properties as a portable, practice-resistant measure of sustained attention \parencite{Dinges1985_pvt,Basner2011_pvt,Loh2004_pvt_short,Basner2018_pvt_practice}.}

\textcolor{blue}{The protocol tests three pre-specified hypotheses. H1 (feasibility): the per-session data yield will remain above 70\% of expected samples on at least three of the four modalities for more than 75\% of recording weeks, confirming that the platform is operationally viable under ordinary teaching conditions. H2 (convergent validity): the participant-session correlation (Spearman's $\rho$) between the within-session $\acute{\text{o}}$ptimo-rate and the PVT change score will be positive and statistically distinguishable from zero across the four-month deployment. H3 (label reliability): Cohen's $\kappa$ between the teacher's prompt-time annotation and a back-coded post-session reconstruction, assessed on a spot-check subset of sessions, will exceed 0.60, indicating adequate within-teacher consistency for a binary judgement. Failure to confirm H1 triggers a platform redesign before any predictive analysis (falsifiability criterion). Failure to confirm H2 or H3 will be reported as an empirical finding about the measurement instruments and will inform the choice of label scheme in subsequent compendium articles. These hypotheses bind the protocol's outcome to observable evidence rather than retrospective interpretation.}

This article documents the protocol. It is observational, prospective and longitudinal, with no educational intervention. Recordings span four months of ordinary teaching in two groups of 3rd and 4th of ESO (Spanish compulsory secondary education, ages 14--16) in a single secondary school, with a comparison-group classroom of the same school year. Four passive data streams are acquired in synchrony: (i) a UCLM ad-hoc smart-classroom platform sampling temperature, relative humidity, CO\textsubscript{2}, PM2.5, PM10, illuminance and sound pressure level at a nominal 15 Hz, with no audio waveform retained; (ii) EmotiBit wristbands recording heart rate, heart-rate variability, electrodermal activity, peripheral skin temperature and \textcolor{blue}{9-axis IMU kinematics} on six rotating units; (iii) Bitbrain Versatile 32-channel semi-dry EEG at 256 Hz on one stratified-sampled participant per week; and (iv) overhead 4K dual-lens video from which only body-pose features are extracted, with the raw footage deleted within 72 hours of feature extraction under a verifiable governance log. The cognitive-performance label is a binary teacher annotation per 5-minute window, anchored to an explicit operational rule (predominantly on-task vs.\ not); external validation is provided by the Psychomotor Vigilance Task administered monthly before and after instrumented sessions, complemented by an end-of-session 0--10 global rating. The protocol has been submitted to the Ethics in Social Research Committee of the Universidad de Castilla-La Mancha (CEIS-UCLM) and operates under a privacy-by-design regime aligned with Article 9 of the General Data Protection Regulation and the Spanish LOPDGDD \parencite{Macenaite2017_minors_gdpr,EDPB2020_consent}.

This study protocol is the methodological backbone of a doctoral compendium. It is observational by design: it does \emph{not} include the active-breaks intervention pursued, in primary education with gifted students, by the parent national project PID2022-137397NB-I00. The two designs share a measurement philosophy but no causal claim is transferable from one to the other. \Cref{sec:methods,sec:results,sec:discussion} document the protocol following the \emph{JMIR Research Protocols} structure, followed by Acknowledgements, Conflicts of Interest, Data Availability and a Multimedia Appendix on the operational privacy-by-design pipeline.
